% =============================================================================
% ISI Country Brief — Austria
% External Supplier Concentration Profile
%
% ISI Paper Series — Country Brief
% International Sovereignty Institute — Research Unit
% March 2026
% =============================================================================
\documentclass[12pt,a4paper]{article}

% --- ISI Paper Series shared template ----------------------------------------
\usepackage{isi-series}

% --- PDF metadata ------------------------------------------------------------
\hypersetup{
  pdftitle={External Supplier Concentration Profile: Austria — ISI EU-27 2024},
  pdfauthor={International Sovereignty Institute -- Research Unit},
  pdfcreator={International Sovereignty Institute},
  pdfsubject={ISI Paper Series -- Country Brief AT, EU-27, Vintage 2024},
  pdfkeywords={sovereignty index, EU-27, supplier concentration,
    austria, country brief},
}

% ==============================================================================
\begin{document}
\hypersetup{pageanchor=false}

% ---------- Title Page with Metadata ------------------------------------------
\begin{titlepage}
\centering
\vspace*{2cm}
{\Large\textsc{ISI Paper Series --- Country Brief}}\\[1.2cm]
{\LARGE\bfseries External Supplier Concentration Profile:\\[0.3em]
Austria}\\[1.5cm]
{\large International Sovereignty Index (ISI)\\[0.3em]
EU-27 --- Vintage 2024 --- Methodology v1.0}\\[1.2cm]
{\large International Sovereignty Institute --- Research Unit}\\[0.8cm]
{\large March 2026}

\vspace{0.8cm}
\noindent\rule{\textwidth}{0.4pt}
\vspace{0.4cm}

% --- Research Metadata --------------------------------------------------------
\begin{flushleft}
\noindent\textbf{Data basis:} ISI Paper~2 --- EU-27 Empirical Results 2024 (Methodology~v1.0, Vintage~2024).\\[4pt]
\noindent\textbf{Methodology:} ISI Paper~1 --- Methodological Foundations, Version~1.0.

\medskip

\noindent\textbf{JEL Codes:} F02, F15, F52, O30, Q43\\[4pt]
\noindent\textbf{Keywords:} sovereignty index; EU-27; supplier concentration; composite indicator; variance decomposition; defence dependence

\medskip

\noindent\textbf{Related publications in the ISI Paper Series:}
\begin{itemize}[nosep,leftmargin=1.5em]
\item ISI Paper~1 --- Technical Specification (v1.0) --- doi:\\
  \url{https://doi.org/10.5281/zenodo.18764227}
\item ISI Paper~2 --- EU-27 Empirical Results 2024 (v1.0) --- doi:\\
  \url{https://doi.org/10.5281/zenodo.18764170}
\end{itemize}
\end{flushleft}
\vfill
\end{titlepage}

\hypersetup{pageanchor=true}
\pagenumbering{arabic}
\pagestyle{fancy}
\fancyhead[L]{\small\sffamily\textit{ISI Country Brief}}
\fancyhead[R]{\small\sffamily\textit{Austria}}
\renewcommand{\headrulewidth}{0.4pt}

% ---------- Body --------------------------------------------------------------

\section{Executive Summary}
\label{sec:summary}

Austria ranks~8th in the EU-27 with a composite \ISI score of~0.383, classified as \emph{moderately concentrated}.  The dominant axes are defence~(0.585) and logistics~(0.526).  The structural profile is characterised as dual-axis concentration (defence and logistics).  Under Dirichlet weight perturbation (10\,000~Monte Carlo draws), Austria's mean rank is~8.04 with a standard deviation of~1.09; the 5th--95th percentile band spans ranks~6 to~10.


\section{Country Position in the EU-27}
\label{sec:position}

Austria occupies rank~8 of~27 EU member states (composite score~0.383).  The EU-27 mean composite score is~0.344 with a standard deviation of~0.070.  Austria's score lies above the EU-27 mean, indicating above-average external supplier concentration.  In the ranking, Austria is situated between Sweden (rank~7; 0.399) and Italy (rank~9; 0.374).


\section{Axis Profile}
\label{sec:axis-profile}

\Cref{tab:axis-profile} presents Austria's scores on all six \ISI axes.

\begin{table}[htbp]
\centering\small
\caption{Austria's \ISI axis profile.}
\label{tab:axis-profile}
\begin{tabular}{@{}lr@{}}
\toprule
\textbf{Axis} & \textbf{Score} \\
\midrule
    Finance           & 0.149 \\
    Energy            & 0.456 \\
    Technology        & 0.261 \\
    Defence           & 0.585 \\
    Critical Inputs   & 0.319 \\
    Logistics         & 0.526 \\
\bottomrule
\end{tabular}
\end{table}

The highest axis score is defence~(0.585), which lies above the EU-27 axis mean of~0.573.  The second-highest axis is logistics~(0.526), above the EU-27 axis mean of~0.518.  The lowest axis score is finance~(0.149), near the EU-27 mean of~0.148.


\section{Structural Drivers}
\label{sec:drivers}

\begin{sloppypar}
Austria's composite score is driven by two axes in combination: defence~(0.585) and logistics~(0.526).  The gap between these two axes is~0.059, indicating a dual-axis concentration profile.  Neither axis alone dominates the composite; rather, the combination of elevated scores on both dimensions determines Austria's position in the ranking.
\end{sloppypar}


\section{Comparative Context}
\label{sec:comparative}

Across the EU-27, the covariance-based variance decomposition shows that defence (35.1\,\%) and logistics (34.3\,\%) together account for 69.4\,\% of total cross-sectional composite variance.  Austria's defence score~(0.585) lies above the EU-27 defence mean~(0.573), and its logistics score~(0.526) lies above the EU-27 logistics mean~(0.518).  Under Dirichlet weight perturbation, Austria's rank standard deviation is~1.09, indicating moderate rank stability.

The composite gap to Austria's immediate neighbours in the ranking is narrow: 0.016 to Sweden~(rank~7) and 0.009 to Italy~(rank~9).  Gaps of this magnitude imply that even modest axis-score shifts could alter ordinal position, a pattern consistent with the moderate rank standard deviation reported above.

Austria's dual-axis profile aligns with the EU-wide variance structure: the two axes that contribute most to cross-sectional composite variance---defence and logistics---are precisely the axes on which Austria scores highest.  Consequently, Austria's position in the ranking is shaped primarily by the dimensions along which member states differ most.

The remaining four axes---finance~(0.149), technology~(0.261), critical inputs~(0.319), and energy~(0.456)---cluster closer to their respective EU-27 means, reinforcing the characterisation of Austria's profile as driven by the defence--logistics pair rather than by broad-based elevation across all dimensions.  The 5th--95th~percentile rank band of~6 to~10 reflects this structure: sufficient concentration to remain in the upper third of the ranking, but not enough separation from neighbouring scores to anchor a narrower band.


\section{Interpretation Boundary}
\label{sec:interpretation}

The \ISI measures concentration of external suppliers, not sovereignty in a normative or political sense.  High concentration may reflect geography, economic specialisation, or alliance structures.  A high composite score does not imply policy failure; a low score does not imply policy success.  The index provides an empirical measurement basis --- what policymakers derive from it is a separate question.


% ---------- References --------------------------------------------------------
\clearpage
\vspace*{1.5cm}

\begingroup\emergencystretch=2em\hbadness=10000
\section*{References}

\begin{enumerate}[label={[\arabic*]},leftmargin=2em,itemsep=6pt]
\item International Sovereignty Institute.
  \textit{International Sovereignty Index (ISI) technical specification, version~1.0}.
  Zenodo, 2026.\\ISI Paper Series, Paper~1.
  doi:\,\href{https://doi.org/10.5281/zenodo.18764227}{10.5281/zenodo.18764227}.

\item International Sovereignty Institute.
  \textit{International Sovereignty Index (ISI): EU-27 empirical results 2024 (v1.0)}.
  Zenodo, 2026.\\ISI Paper Series, Paper~2.
  doi:\,\href{https://doi.org/10.5281/zenodo.18764170}{10.5281/zenodo.18764170}.
\end{enumerate}
\endgroup

\end{document}
